%==================================================================
\section{Plan de calcul et reproduction (23 septembre 2026)}\label{sec:plan}
\emph{Pour la personne qui fournira le calcul et pour celle qui vérifiera. Chiffres mesurés sur les exécutions de la nuit (fichiers
\code{cubes/*/results.txt}, \code{NUIT\_JOURNAL.md}).}

\subsection{Ce que décide le calcul complet}
Chaque étape est exhaustive : graphes $F$ (26 sur 156), classes de phases (46), multi-ensembles de supports par classe, puis étage quotient
découpé en cubes qui couvrent tous les cas (chaque orbite indépendante a exactement deux doubles). Si tout va au bout :
\begin{itemize}[nosep]
\item \textbf{tout est UNSAT} : aucun $\srg(99,14,1,2)$ n'a d'automorphisme d'ordre 3 (sous réserve des points de la section suivante) ;
\item \textbf{un problème quotient est SAT} : ce n'est \emph{pas} encore une réponse, car l'étage quotient ignore les phases du bloc $C$.
Il faut alors l'étage des phases (\code{stage\_p.py}, prêt et testé sur le graphe tour) : soit on obtient un graphe, vérifiable en une seconde
($A^2+A=12I+2J$), soit ce $Q$ est exclu et il faut traiter les autres solutions quotient de l'instance. Ce cas n'est pas chiffré.
\end{itemize}
{\sloppy Indices à ce jour : 10 problèmes quotient entièrement UNSAT, 3\,381 cubes UNSAT à l'arrêt (\code{MANIFEST.txt}), \textbf{aucune solution nulle part}, aucune solution
invariante pour l'instance la plus symétrique. Cela rend UNSAT très probable, sans le démontrer.\par}

\subsection{Niveau de confiance de chaque maillon}
\begin{center}\small\begin{tabular}{p{0.28\textwidth}p{0.30\textwidth}p{0.34\textwidth}}
maillon & garanti par & ce qui manque\\\hline
absence de point fixe de $g$ & articles de revues (Makhnev--Minakova 2004, Behbahani--Lam 2011) & --\\
$\tau=6$, graphes $F$, classes & démonstrations (dont une preuve autonome de $\tau=6$), certificats de Farkas & relecture humaine\\
dépôts GitHub de tiers & -- & aucun résultat n'en dépend (contexte seulement)\\
lemmes des entrées, LQ, niveau 2.5 & démonstrations (\S\ref{sec:l25}, \S\ref{sec:nuit}) & \textbf{relecture humaine du niveau 2.5}\\
cassage de symétrie, couverture des cubes & lemme~\ref{lem:symlex}, validations exhaustives & relecture humaine\\
encodages CNF & validations sur graphes réels et témoins SAT & (preuve formelle hors de portée)\\
UNSAT de l'étage quotient & DRAT + drat-trim (9 instances) & DRAT cube par cube pour les autres\\
\textbf{complétude des énumérations} & « infaisable » de CP-SAT ; certificat DRAT (\code{cert\_enum.py}) & certifiée pour 7 classes
(0, 4, 15, 25, 30, 32, 42) ; \textbf{35 autres complètes à certifier ; 29, 43, 44, 45 incomplètes}
\end{tabular}\end{center}

\paragraph{Certifier une énumération.} L'« infaisable » final de CP-SAT n'est pas vérifiable : sur un modèle test, CP-SAT conclut dès le
pré-traitement (symétries, réductions duales) et n'émet aucune preuve, bien qu'il en expose l'option. \code{cert\_enum.py} réécrit donc le modèle
d'énumération avec phases entièrement en CNF, contrainte par contrainte : supports triés, $G_0$, TY, sommes de lignes, doubles, L1, C3, colonne
phasée par orbite (parmi celles qu'autorise le filtre de type), égalités résiduelles, condition des triplets. On y ajoute la clause bloquante de
chaque multi-ensemble trouvé ; l'insatisfiabilité, prouvée par CaDiCaL et vérifiée par drat-trim, certifie que la liste est complète. Une
contrainte de symétrie est ajoutée : deux orbites de même support ont des indices de colonne rangés dans l'ordre. Elle est valide, car échanger
deux telles orbites (bloc $C$ compris) transforme une solution en solution.
\emph{Validation dans les deux sens.} (i) Toute solution de la CNF est relue par un évaluateur Python indépendant qui teste chaque contrainte du
modèle ; (ii) chaque multi-ensemble trouvé par CP-SAT, débloqué seul, redonne SAT avec exactement ce multi-ensemble. Classe 15 (4,6\,M clauses) :
débloquée SAT (0 violation), bloquée UNSAT. Classe 0 : chacun des 4 multi-ensembles débloqué seul redonne SAT avec lui-même (0 violation), les 4
bloqués donnent UNSAT.
\emph{Certificats} (CaDiCaL 3.0.1, preuve DRAT binaire, drat-trim VERIFIED) :
\begin{center}\footnotesize\begin{tabular}{rrlrrr}
classe & multi-ens. & CNF (sha256, début) & résol. & preuve & vérif.\\\hline
0 & 4 & \code{f791a54edc40d441} & 202\,s & 27\,Mo & 174\,s\\
4 & 96 & \code{32d5e21000cfa5d5} & 120\,s & 23\,Mo & 38\,s\\
15 & 1 & \code{6e4274b135ed5b96} & 272\,s & 30\,Mo & 224\,s\\
25 & 1 & \code{4a955e8e9e58b6b9} & 1\,179\,s & 84\,Mo & 330\,s\\
30 & 1 & \code{c2e9662631bbf6cc} & 140\,s & 45\,Mo & 67\,s\\
32 & 1 & \code{4cfb3bd63aa6c23d} & 180\,s & 61\,Mo & 79\,s\\
42 & 1 & \code{b096ddbee272e204} & 863\,s & 76\,Mo & 611\,s
\end{tabular}\end{center}
Le coût est de quelques minutes à une demi-heure par classe. Réénumération avec phases : 42 classes sur 46 ont une liste complète (un seul
multi-ensemble compatible par classe, sauf les classes 0 et 27, qui en ont 4, et la classe 4, qui en a 96) ; restent incomplètes 29 (31 trouvés),
43 (2 trouvés), 44 et 45. Les 35 listes complètes non encore certifiées le seront par \code{cert\_enum\_all.sh} (image \code{c3box}), environ
6 à 10\,h sur une machine.

\subsection{Coût mesuré et estimation}
Avec la méthode incrémentale, les passes à 5\,000 puis 50\,000 conflits règlent 85 à 95\,\% des 325 cubes d'une instance ; un cube dur coûte
en moyenne 1 à 45\,min selon l'instance (moyennes mesurées : 1\,min pour m00028\_sL0, 15 pour m16420\_sL1, 22 pour m12295\_s0015, 45 pour
m23431\_s0028). m00028\_sL0 a été réglée entièrement en environ 3\,h de CPU.
\begin{center}\small\begin{tabular}{lr}
poste & estimation (heures-CPU)\\\hline
11 instances en cours (environ 200 cubes durs restants) & 50--100\\
30 instances non commencées (dont les 18 de la classe 4), 3--6\,h chacune & 90--180\\
certification DRAT de tous les UNSAT & du même ordre que la résolution\\
classes 29, 43, 44, 45 & \textbf{inconnu} (principal risque)
\end{tabular}\end{center}

\subsection{Passage à l'échelle}
\begin{itemize}[nosep]
\item Environ 40 instances $\times$ 325 cubes, soit environ 13\,000 tâches indépendantes, plus les sous-cubes (25 par nœud dur, eux aussi
indépendants). Chaque tâche : 1 cœur, environ 300\,Mo. Machine de calcul du projet (16 threads, 32\,Go) : 12 à 14 runners, un à deux jours pour la
partie chiffrée. 100 cœurs : 2 à 3\,h. 1\,000 cœurs : moins d'une heure, la durée étant alors bornée par le cube le plus long (à redécouper).
\item Limite : l'énumération CP-SAT d'une classe n'utilise bien que 8 à 16 threads ; pour des centaines de cœurs, il faut la découper elle aussi
(par exemple selon le support de la première orbite), à concevoir.
\item Les clusters n'acceptent généralement pas Docker : Apptainer/Singularity importe les archives d'images telles quelles.
\end{itemize}

\subsection{Reproduction octet pour octet}
{\sloppy
Trois niveaux, qu'il faut distinguer :
\begin{enumerate}[nosep]
\item \textbf{Identique à l'octet, vérifié par SHA-256.} Fichiers d'instances ; CNF (vérifié : m16420\_sL0 donne le même SHA-256
\code{a706468daf92ee33\ldots} avec deux graines de hachage Python différentes et dans trois générations indépendantes) ; images Docker (archives
exportées, \code{docker/IMAGES.sha256}, et identifiants d'image à retrouver après \code{docker load}, \code{docker/IMAGES.txt}) ; résultats sous
forme canonique triée (\code{manifest.py}, \code{MANIFEST.txt}). Vérification croisée : les listes de supports calculées par trois méthodes
différentes (LQ seule, réduction par $\operatorname{Aut}(F)$, phases) ont des SHA-256 canoniques identiques sur les masques 28, 8290, 9313, 13452
et 29448.
\item \textbf{Déterministe si l'on configure à l'identique.} CaDiCaL sur un thread, avec des budgets de \emph{conflits} et non de secondes, redonne
le même verdict, et même la même preuve DRAT octet pour octet. Les exécutions de cette nuit ont utilisé des limites de temps et ont été
interrompues puis reprises : leurs journaux ne se rejouent pas à l'identique, leurs verdicts si.
\item \textbf{Non reproductible par nature.} CP-SAT à plusieurs threads (ordre de découverte, durées). Seul l'\emph{ensemble} final compte, d'où
la forme canonique.
\end{enumerate}
Critère scientifique retenu, plus fort qu'un simple rejeu : \textbf{même CNF (SHA-256) et certificat DRAT vérifié par drat-trim}. Ce contrôle ne
dépend ni du solveur, ni de la machine, ni de l'ordre d'exécution.\par}

\subsection{Ordre des opérations recommandé}
{\sloppy
\begin{enumerate}[nosep]
\item Charger les images : \code{docker load < docker/c99audit.image.tar.gz} (idem \code{proofbox}, \code{texbox}) ; comparer les
identifiants avec \code{docker/IMAGES.txt}. À défaut : \code{docker/Dockerfile.c99audit}, \code{docker/proofbox/Dockerfile} (versions figées,
mais binaires pas garantis identiques).
\item Lancer les runners de cubes (un par cœur) :
\code{docker run -d --cpus 1 -m 1500m -v "\$PWD:/w" -w /w c99audit sh cube\_runner\_inc.sh 5000} ; la file est \code{cubes/LIST.txt}.
\item En parallèle, les énumérations des classes 29, 43, 44, 45 :
\code{python enum\_supports\_phased.py <classe> 3600 1000 8}, puis \code{make\_phased\_instances.py} et ajout à la file.
\item Certifier : \code{cert\_sym.sh} pour une instance entière ; pour un cube, CNF de base plus ses clauses unitaires, même vérification.
\item Certifier les énumérations : \code{sh cert\_enum\_all.sh} dans l'image \code{c3box} (génère chaque CNF, puis CaDiCaL, preuve DRAT et
drat-trim, une classe à la fois ; résultats dans \code{certs/SUMMARY\_ENUM.txt}).
\item Régénérer \code{MANIFEST.txt} et le comparer ligne à ligne avec toute autre reproduction.
\end{enumerate}\par}
